\documentclass[11pt,a4paper]{article}
\usepackage[lmargin=1.0in,rmargin=1.0in,bottom=1.0in,top=1.0in,twoside=False]{geometry}
\usepackage{mathtools}
\usepackage{hchang}

\let\proof\relax
\let\endproof\relax
\let\qedhere\relax

\usepackage{amsthm}
%\usepackage{amsmath}
\usepackage{graphicx}%,algorithmic,algorithm}%, verbatim}
\usepackage{enumerate}
\usepackage[scale=0.88]{FiraMono}
\usepackage{xcolor}
\usepackage{listings}
\usepackage{tikz}
\usepackage{stackengine}
\usepackage{accents}
\usepackage{textpos}
\usetikzlibrary{shapes,calc,fit,arrows.meta,positioning}

\usepackage[T1]{fontenc}

\usepackage{enumitem}

\usepackage{microtype}
%\usepackage{amsfonts}
\usepackage{comment}
\usepackage[capitalize,nameinlink]{cleveref}

\renewcommand{\figurename}{Figure}
\crefname{figure}{Figure}{Figures}
\Crefname{figure}{Figure}{Figures}

\usepackage{xcolor}
\usepackage{pifont}

% LAX website surfaces.
\definecolor{LaxCodeBg}{HTML}{F6F8FB}
\definecolor{LaxBorder}{HTML}{E8E3DC}
\definecolor{LaxText}{HTML}{24292E}
\definecolor{LaxDim}{HTML}{6B7280}
\definecolor{LaxLineNo}{HTML}{B3B8C0}
\definecolor{LaxAccent}{HTML}{3B6B9A}
\definecolor{LaxProvenBg}{HTML}{E8F5E8}

% Shiki "github-light" token colours used by laxarchive.org.
\definecolor{GHKeyword}{HTML}{D73A49}
\definecolor{GHConstant}{HTML}{005CC5}
\definecolor{GHEntity}{HTML}{6F42C1}
\definecolor{GHString}{HTML}{032F62}
\definecolor{GHComment}{HTML}{6A737D}
\definecolor{GHVariable}{HTML}{E36209}

\lstdefinelanguage{LeanFour}{
  sensitive=true,
  alsoletter={_,'},
  morecomment=[l]{--},
  morecomment=[s]{/-}{-/},
  morestring=[b]",
  morekeywords=[1]{import,namespace,end,def,axiom,theorem,by,open,only},
  morekeywords=[2]{Nat,Prop,rfl},
  morekeywords=[3]{Mathlib,LaxDemo,EvenAddition,LaxDemoProofs,twice,Even,even_add},
  morekeywords=[4]{rintro,refine,simp,ac_rfl},
  morekeywords=[5]{m,n,a,b,k},
  keywordstyle=[1]\color{GHKeyword},
  keywordstyle=[2]\color{GHConstant},
  keywordstyle=[3]\color{GHEntity},
  keywordstyle=[4]\color{GHEntity},
  keywordstyle=[5]\color{GHVariable},
  commentstyle=\color{GHComment},
  stringstyle=\color{GHString},
  literate=
    {→}{{$\to$}}1
    {∃}{{$\exists$}}1
    {⟨}{{$\langle$}}1
    {⟩}{{$\rangle$}}1
}

\lstdefinestyle{LaxLean}{
  language=LeanFour,
  basicstyle=\ttfamily\fontsize{7.5}{11.9}\selectfont\color{LaxText},
  numbers=left,
  numberstyle=\ttfamily\fontsize{6.4}{11.9}\selectfont\color{LaxLineNo},
  numbersep=7pt,
  stepnumber=1,
  xleftmargin=22pt,
  xrightmargin=7pt,
  aboveskip=0pt,
  belowskip=0pt,
  columns=fullflexible,
  keepspaces=true,
  showstringspaces=false,
  tabsize=2,
  breaklines=false,
  frame=none,
  backgroundcolor=\color{LaxCodeBg},
  linebackgroundheight=8.1pt,
  linebackgrounddepth=3.8pt
}

\newcommand{\ConceptLineBackground}{%
  \color{LaxCodeBg}%
  \ifnum\value{lstnumber}>9\relax
    \ifnum\value{lstnumber}<13\relax\color{LaxProvenBg}\fi
  \fi
}
\newcommand{\ProofLineBackground}{%
  \color{LaxCodeBg}%
  \ifnum\value{lstnumber}>6\relax
    \ifnum\value{lstnumber}<14\relax\color{LaxProvenBg}\fi
  \fi
}

\newlength{\LaxPanelWidth}
\newlength{\LaxPanelHeight}
\setlength{\LaxPanelWidth}{7.45cm}
\setlength{\LaxPanelHeight}{7.75cm}

\newcommand{\LaxHeader}[2]{%
  \noindent\hspace*{3.1mm}%
  {\footnotesize\scshape #1}%
  \hfill
  {\ttfamily\scriptsize\color{LaxDim}#2}%
  \hspace*{3.1mm}\par
  \vspace{1.6mm}%
  {\color{LaxBorder}\hrule height 0.45pt}%
  \vspace{1.7mm}%
}

\newsavebox{\ConceptBox}
\newsavebox{\ProofBox}


% Exact colours from laxarchive.org
\definecolor{LaxThmBackground}{HTML}{E8F5E8}
\definecolor{LaxThmForeground}{HTML}{2D6A2D}
\definecolor{LaxThmBorder}{HTML}{B8DAB8}

\DeclareRobustCommand{\laxthm}{%
  \tikz[baseline=(LaxTHM.base)]{%
    \node[
      draw=LaxThmBorder,
      fill=LaxThmBackground,
      text=LaxThmForeground,
      line width=.75bp,       % 1 CSS pixel
      rounded corners=1.5bp, % 2 CSS pixels
      inner xsep=3bp,         % 4 CSS pixels
      inner ysep=.75bp,       % 1 CSS pixel
      outer sep=0pt,
      font=\fontsize{7.65bp}{9.18bp}\selectfont
            \rmfamily\bfseries
    ] (LaxTHM) {%
      T\kern.04em
      H\kern.04em
      M\kern.04em
      \raisebox{-.08ex}{\ding{51}}%
    };%
  }%
  \kern1.5bp%                
}

\usepackage{stmaryrd}

\renewcommand{\ge}{\geqslant}
\renewcommand{\geq}{\geqslant}
\renewcommand{\le}{\leqslant}
\renewcommand{\leq}{\leqslant}

\newtheorem{lemma}{Lemma}
\newtheorem{theorem}{Theorem}
\newtheorem{question}{Question}
\newtheorem{corollary}{Corollary}
\newtheorem{definition}{Definition}
\newtheorem{conjecture}{Conjecture}
\newtheorem{observation}{Observation}
\newtheorem{claim}{Claim}
\newcommand{\Oh}{\mathcal{O}}
\newcommand{\dist}{\mathsf{dist}}
\newcommand{\prof}[2]{\mathsf{prof}_{#1}[#2]}
\newcommand{\bover}[1]{\accentset{\bullet}{#1}}
\newcommand{\real}{\mathbb{R}}


\title{Lax}

\begin{document}

\date{}

\maketitle

\begin{abstract}
  We introduce lax, an open archive of formalized mathematics.
  This document outlines what is lax, how you can and why you should contribute to it, and what guarantees it provides. 
\end{abstract}

\section{What is lax?}

Lax is an archive of formalized mathematics, which comprises:
\begin{itemize}
\item a command-line interface \href{https://www.npmjs.com/package/lax-archive}{\textsf{lax-CLI}}, which primarily allows one to submit new entries,
\item a database \href{https://github.com/lax-archive/lax-database/}{\textsf{lax-archive}}, which records the metadata of every submission, 
\item a \href{https://laxarchive.org/}{website}, which enables to browse these submissions.
\end{itemize}
The entire project is fully open: the source code of \textsf{lax-CLI} and of the website, the product specification, the database content, the submissions, and the generated artifacts are all public and reusable.
Currently, \href{https://lean-lang.org/}{Lean} is the only supported language for the submissions, the `l' of lax stands for Lean, and `ax' for \textbf{a}r\textbf{ch}ive.

\medskip

Unlike a~library, lax targets contemporary mathematical results, which would take considerable time and effort to be manually formalized.
Typically lax proofs are written by AI agents.
They are not held to any beauty cannons, as long as Lean's kernel certifies that they indeed prove what the submission claims. 
Lax does however care that these claims (the \emph{concepts}, see next section), which constitute the exposed surface on the website, are transparently written and can easily be evaluated and reviewed by humans.

\section{What is a lax submission?}

A lax submission is a~commit on a~public Git repository which contains appropriate files and folders.
Rather than going through the whole specification, which can be consulted by humans and AI agents (for instance, by running `lax spec'), we highlight its salient features.

A lax submission separates in two different folders \emph{concepts} and \emph{proofs}.
In the concept folder, each Lean file has a~possibly empty list of definitions\footnote{introduced with the Lean keywords \textsf{def}, \textsf{structure}, \textsf{inductive}, \textsf{abbrev}, etc.} and at most one statement introduced as axiom.
The content of the concept folder is semantically closed: the meaning of every declaration (definition or axiom statement) can be evaluated within the union of the concept folder and Mathblib. 
In the proof folder, Lean files are unconstrained.
In a~typical submission, every axiom declared in the concept files is restated in the proof folder as a~theorem and proven (under the basic Lean axioms propext, Quot.sound, Classical.choice).

The content of concept files is exposed on the website.
Ideally, the declarations are transparently translated from their natural-language counterparts, and can be understood and reviewed by users without a~strong Lean or proof-assistant background.
In particular, Lean tactics are unwelcome in concept files.
Every axiom statement accompanied by a~proof (found in the proof package) is marked as proven (green background or \laxthm) on the website.

\begin{figure}[!ht]
  \centering
  \begin{lrbox}{\ConceptBox}
\begin{minipage}[t][\LaxPanelHeight][t]{\LaxPanelWidth}
  \vspace*{2.5mm}
  \LaxHeader{Concept}{LaxDemo/EvenAddition.lean}
\begin{lstlisting}[style=LaxLean,linebackgroundcolor=\ConceptLineBackground]
import Mathlib

namespace LaxDemo.EvenAddition

def twice (n : Nat) : Nat := n + n

def Even (n : Nat) : Prop :=
  ∃ k, n = twice k

axiom even_add {m n : Nat} :
    Even m → Even n →
    Even (m + n)

end LaxDemo.EvenAddition
\end{lstlisting}
\end{minipage}
\end{lrbox}

\begin{lrbox}{\ProofBox}
\begin{minipage}[t][\LaxPanelHeight][t]{\LaxPanelWidth}
  \vspace*{2.5mm}
  \LaxHeader{Proof}{LaxDemoProofs/EvenAddition.lean}
\begin{lstlisting}[style=LaxLean,linebackgroundcolor=\ProofLineBackground]
import LaxDemo.EvenAddition

namespace LaxDemoProofs

open LaxDemo.EvenAddition

theorem even_add {m n : Nat} :
    Even m → Even n →
    Even (m + n) := by
  rintro ⟨a, rfl⟩ ⟨b, rfl⟩
  refine ⟨a + b, ?_⟩
  simp only [twice]
  ac_rfl

end LaxDemoProofs
\end{lstlisting}
\end{minipage}
\end{lrbox}


\begin{tikzpicture}[
  panel/.style={
    draw=LaxBorder,
    fill=LaxCodeBg,
    rounded corners=2mm,
    line width=0.55pt,
    inner sep=0pt,
    outer sep=0pt,
    anchor=north west
  }
]
  \node[panel] (concept) {\usebox{\ConceptBox}};
  \node[panel, right=17mm of concept.north east, anchor=north west] (proof)
    {\usebox{\ProofBox}};

  \draw[-{Stealth[length=2.2mm,width=1.5mm]},
        draw=LaxAccent, line width=0.75pt]
    ([xshift=2.0mm,yshift=-5.25cm]concept.north east) --
    node[midway, above=1.6mm, fill=white, inner xsep=1.0mm, inner ysep=0.4mm,
         text=LaxDim, font=\scriptsize\itshape] {proved by}
    ([xshift=-2.0mm,yshift=-5.25cm]proof.north west);
\end{tikzpicture}


  \caption{Left: A concept file, as displayed on the website. Right: A matching proof file, as found in the submitter's repository.}
  \label{fig:concept-proof}
\end{figure}

By default, a~lax submission initially is a~\emph{draft}.
In this state, it can be modified ad libitum by resubmitting.
A submission can be made \emph{permanent}, and then becomes immutable.
Lax submissions can reuse and build on past permanent entries.

\section{Why using lax?}

The starting point is that you already are convinced with the benefits of having results you care about formalized.
Year 2026 was when frontier AI agents unlocked the ability of fully autonomously translating natural-language proofs, however long they were, into Lean.
Consequently, at a~small (and ever smaller) cost, and at the time of telling your agent which result you want formalized next, you very likely find your wish granted when you later come back to the session.   

Once you thus get a~Lean proof of, say, your latest result, why should you further make the effort of submitting it to lax?
Why? Because, this way, your work
\begin{itemize}
\item is more visible: now a~searchable platform (attracting a~growing community) points to it.
\item is more legible: the presentation work is automated away by the submission guidelines, the submitting agent, and the website.
  The suitable subset of the work that should be exposed is conveniently displayed, with semantic hyperlinks, proof DAGs, concept dependencies, etc.    
\item can be more easily built upon: \textsf{lax-CLI} makes it simple to fetch results from past submissions, in view of completing new Lean projects. 
\item will be validated: users browsing your concept files can vet that those are formalizing the intended mathematical definitions and statements, or otherwise, raise concerns.  
\end{itemize}
Again, AI agents can autonomously and reliably follow \textsf{lax-CLI}'s documentation to prepare lax submissions, monitor ongoing submissions, and troubleshoot if need be.
The human in the loop mainly checks that the abstract, commented concept files, and optional implementation notes are clean and as intended.
This does not take much time.
As the friction quickly evaporates, do not keep your formalized results isolated.

\section{Can you trust lax?}

If the axiom of a~concept appears as proven on the website, should I trust that the corresponding proof package actually proves it?
Assuming you already believe in the soundness of Lean's kernel, you have to further trust us on basically the entire submission pipeline.
This is too much to ask, even with everything open-source.

Let us summarize what we claim this pipeline does.

Of course, the Lean code of every submission is publicly available.
So for any given submission, you can independently check that 
  
\section{Can lax trust you?}

Currently it does.
The community is still small enough that imagining adversarial submissions or attacks.

\section{How is lax built?}

The current lax infrastructure relies on GitHub services:
\begin{itemize}
  \item the project repositories are hosted on GitHub,  
  \item the submission pipeline uses GitHub Actions,
  \item the website is built and deployed via GitHub Pages, and
  \item \textsf{lax-CLI} is available on npm.
\end{itemize}

\bibliographystyle{plain}
\bibliography{main}

\end{document}
